% Standalone research entry 211; batch 208--217.
% Original section material preserved from Pasted text(3).txt.
% Research additions: 27 September 2026.
% Compile twice with: lualatex 211.tex
% !TeX program = lualatex
% Compile twice: lualatex open_problems_500.tex
% All references and prose are embedded; no BibTeX database or external inputs.
\documentclass[11pt,a4paper]{article}
\usepackage[margin=24mm,headheight=14pt]{geometry}
\usepackage{fontspec}
\setmainfont{Latin Modern Roman}
\setsansfont{Latin Modern Sans}
\setmonofont{Latin Modern Mono}
\usepackage{amsmath,amssymb,mathtools}
\usepackage{unicode-math}
\setmathfont{Latin Modern Math}
\usepackage{microtype}
\usepackage{enumitem,xurl,needspace}
\usepackage[dvipsnames]{xcolor}
\usepackage{hyperref}
\hypersetup{unicode=true,colorlinks=true,linkcolor=MidnightBlue,urlcolor=MidnightBlue,
 pdftitle={500 Mathematical Problem Statements},pdfauthor={Source-annotated compilation},bookmarksnumbered=true}
\usepackage{bookmark}
\usepackage{tocloft}
\setlength{\cftsecnumwidth}{3em}
\cftsetpnumwidth{2.5em}
\cftsetrmarg{4em}
\usepackage{titlesec}
\titleformat{\section}{\Large\bfseries\raggedright}{\thesection}{0.7em}{}
\usepackage{fancyhdr}
\pagestyle{fancy}\fancyhf{}
\fancyhead[L]{\small\sffamily Mathematical problem statements}
\fancyhead[R]{\small Source edition 22}
\fancyfoot[C]{\thepage}
\setlength{\parindent}{0pt}
\setlength{\parskip}{5pt plus 1pt}
\setlength{\emergencystretch}{3em}
\setcounter{tocdepth}{1}
\setcounter{secnumdepth}{1}
\setlist[itemize]{leftmargin=3em,itemsep=5pt,topsep=4pt,parsep=0pt}
\allowdisplaybreaks
\newcommand{\N}{\mathbb{N}}
\newcommand{\Z}{\mathbb{Z}}
\newcommand{\Q}{\mathbb{Q}}
\newcommand{\R}{\mathbb{R}}
\newcommand{\C}{\mathbb{C}}
\newcommand{\F}{\mathbb{F}}
\newcommand{\A}{\mathbb{A}}
\newcommand{\PP}{\mathbb P}
\newcommand{\E}{\mathbb{E}}
\newcommand{\eps}{\varepsilon}
\newcommand{\dd}{\,\mathrm d}
\newcommand{\sref}[2]{\hyperlink{source:#1:#2}{[S#2]}}
\newcommand{\eref}[2]{\hyperlink{extra:#1:#2}{[E#2]}}
% Turn the next switch off for a shorter reading copy. The quotations preserve
% every exactTarget field, including qualifications omitted from a short summary.
\newif\ifshowcatalogue
\showcataloguetrue
\newcommand{\cataloguescope}[1]{\ifshowcatalogue
 \paragraph{Exact scope in the supplied catalogue.}
 \begin{quote}\small #1\end{quote}\fi}

% Standalone wrapper; no external inputs or bibliography files are required.
\fancyhead[L]{\small\sffamily Research notebook: section 211}
\fancyhead[R]{\small 27 September 2026}
\hypersetup{pdftitle={Research attempt 211},pdfauthor={Research notebook}}
\setlength{\parskip}{4pt plus 1pt}
\setlist[itemize]{before=\small\interlinepenalty10000,itemsep=3pt}
\begin{document}
\setcounter{section}{210}
% ================================================================
% problemId: problem.hopf-conjecture-for-positive-sectional-curvature
% source releaseRank: 211; source releaseStatus: open_with_solved_subcases
% exactTarget SHA-256: 41dd4da598b0e82a42aaa4b8feb9909a3c7306eead7e9e8421c93dda6d228e9e
\clearpage
\section{\texorpdfstring{Hopf conjecture for positive sectional curvature}{Hopf conjecture for positive sectional curvature}}\label{211}
\subsection{Definitions and mathematical statement}
Let $(M^{2m},g)$ be a closed Riemannian manifold. For a two-plane $\sigma=\operatorname{span}(u,v)\subset T_xM$, sectional curvature is
$K_g(\sigma)=\langle R(u,v)v,u\rangle/(|u|^2|v|^2-\langle u,v\rangle^2)$, with the convention making the round sphere positive. The Euler characteristic is
$\chi(M)=\sum_j(-1)^j\dim H_j(M,\Q)$. The conjecture states
\[
 K_g(\sigma)>0\text{ for every }x,\sigma
 \quad\Longrightarrow\quad\chi(M)>0.
\]
No symmetry or simple-connectedness hypothesis is imposed. Positivity of Ricci curvature or scalar curvature alone is a different and weaker curvature condition.
\par\smallskip\textit{Formulation and concept references:} \sref{211}{1}.
\cataloguescope{Prove that every closed even-dimensional Riemannian manifold whose sectional curvatures are everywhere positive has positive Euler characteristic.}
\subsection{Short English statement}
Must every closed even-dimensional manifold with everywhere positive sectional curvature have positive Euler characteristic?
% BEGIN RESEARCH ADDITION 211
\subsection{Research attempt}

\paragraph{Current frontier.}
The unrestricted assertion must be separated from the substantial symmetry
results. Nienhaus proves the positive-sectional-curvature conclusion with an
additional effective isometric $T^4$-action \eref{211}{1}. The 2025 source
\sref{211}{1}, Theorem A and its first corollary, instead treats positive second
intermediate Ricci curvature: dimension divisible by four and $T^{10}$ symmetry
suffice for positive Euler characteristic. Neither result supplies the missing
symmetry in this entry. The calculations below use the classical
Weitzenb\"ock identity, in the normalization where the unit sphere has positive
curvature \eref{211}{2}, Section 3, (3.10). The proposed extensions, and their
failure tests, are this notebook's deductions; no novelty of the general
Bochner method is asserted.

\paragraph{Attack route.}
One sufficient route is to annihilate every odd-degree harmonic form. A second
is to prove a positive integral Euler-density formula without pointwise
positivity. The former makes a particularly testable claim: does positive
sectional curvature force a usable lower bound for the curvature endomorphism
on odd forms? I first test the pointwise version exactly, then replace it by
an integral small-negative-part condition. A negative eigenvalue of an
algebraic curvature term is not a counterexample manifold, but it would rule
out a commonly tempting shortcut before global arguments are attempted.

\paragraph{Attempt.}
\textit{1. The topological target of the analytic argument.}
It is enough to treat each connected component, and then an oriented double
cover if needed: Euler characteristic multiplies by the covering degree.
For a closed connected oriented even-dimensional manifold,
\[
 H^{\mathrm{odd}}(M;\R)=0\quad\Longrightarrow\quad
 \chi(M)=\sum_j b_{2j}(M)\ge2.
\]
This is a sufficient condition, not an asserted equivalent formulation of
Hopf's conjecture. Positive sectional curvature implies positive Ricci
curvature. The identity
\[
 0=\|\nabla\omega\|_2^2+
   \int_M\langle q_p(R)\omega,\omega\rangle\,\mathrm{dvol},
 \qquad \Delta\omega=0,
\]
with $q_1(R)=\operatorname{Ric}$, gives $b_1=0$; Poincar\'e duality also gives
$b_{n-1}=0$ \eref{211}{2}. Thus in dimension four
$\chi=2+b_2>0$. In dimension six the same reasoning leaves
$\chi=2+2b_2-b_3$, displaying exactly the uncontrolled odd degree.

\textit{2. An exact obstruction in dimension six.}
On oriented $\R^4$, write $\Lambda^2=\Lambda^2_+\oplus\Lambda^2_-$ in the
orthonormal bases
\[
 \frac{12+34}{\sqrt2},\ \frac{13-24}{\sqrt2},\
 \frac{14+23}{\sqrt2};\qquad
 \frac{12-34}{\sqrt2},\ \frac{13+24}{\sqrt2},\
 \frac{14-23}{\sqrt2},
\]
where $ij=e_i\wedge e_j$. Consider the symmetric curvature operator
\[
 R_4=\begin{pmatrix}A&0\\0&C\end{pmatrix},\qquad
 A=\operatorname{diag}(-1/2,-1/2,4),\quad C=I_3.
\]
It satisfies the first Bianchi identity because $\operatorname{tr}A=
\operatorname{tr}C=3$. This can also be checked directly: in the lexicographic
wedge basis its only potentially nonzero Bianchi expression with four
distinct indices is
\[
 (R_4)_{12,34}-(R_4)_{13,24}+(R_4)_{14,23}
 =-3/4-3/4+3/2=0.
\]
For a unit decomposable two-form $\eta$, the identity
$\eta\wedge\eta=0$ implies $\|\eta_+\|^2=\|\eta_-\|^2=1/2$.
Consequently
\[
 \langle R_4\eta,\eta\rangle\ge\tfrac12(-\tfrac12)+\tfrac12=\tfrac14.
\]
All sectional curvatures are positive, although the curvature operator itself
has two negative eigenvalues.

Here $\operatorname{Ric}=3I$. On two-forms the Weitzenb\"ock endomorphism is
$6I-2R_4$, so
\[
 \alpha=(14+23)/\sqrt2,\qquad q_2(R_4)\alpha=-2\alpha.
\]
Extend this algebraic tensor by zero to $\R^4\oplus\R^2$ and put
\[
 R_6=R_4\oplus0+\tfrac1{10}I_{\Lambda^2\R^6}.
\]
Projection of a decomposable two-form onto $\Lambda^2\R^4$ is still
decomposable. The extended first term is therefore nonnegative on planes,
and $R_6$ has sectional curvature at least $1/10$ on every plane. On
$\psi=\alpha\wedge e_5$, however, the product representation and the constant
curvature identity $q_p(\kappa I)=\kappa p(n-p)I$ give
\[
 q_3(R_6)\psi=(-2+9/10)\psi=-\tfrac{11}{10}\psi.                 \tag{211.1}
\]
These normalizations can be verified without a curvature-sign convention
imported from another text: if $L_{ij}$ is the infinitesimal orthogonal
rotation acting on exterior forms, set
\[
 q_p(R)=-\sum_{i<j,\,k<\ell}R_{ij,k\ell}L_{ij}L_{k\ell}.
\]
Then $q_p(I)=p(n-p)I$. The accompanying exact rational computation constructs
this $20$-by-$20$ matrix for (211.1). Its eigenvalues, with multiplicities, are
\[
 -11/10\ (2),\quad39/10\ (8),\quad49/10\ (6),\quad79/10\ (4).
\]
The displayed eigenpair and Bianchi identity are checked symbolically, not
inferred from floating-point eigenvalues. Thus the implication
``positive sectional curvature gives positive $q_p$ in every odd degree''
is false already for an algebraic tensor in dimension six.

\textit{3. A replacement allowing a negative part.}
Let $n\ge3$, $V=\operatorname{vol}M$, and suppose a scalar Sobolev inequality
has been established with a constant $S$:
\[
 \|u\|_{2n/(n-2)}^2\le S\bigl(\|du\|_2^2+V^{-2/n}\|u\|_2^2\bigr).
                                                               \tag{211.2}
\]
Fix a degree $p$ and suppose there are $a_p>0$ and a nonnegative function
$W_p\in L^{n/2}$ with
\[
 q_p(R)(x)\ge(a_p-W_p(x))I.
\]
Define $b_p=S\|W_p\|_{n/2}$. If
\[
 b_p<1,\qquad b_p V^{-2/n}<a_p,                                 \tag{211.3}
\]
then $H^p(M;\R)=0$.

Indeed, for a harmonic $p$-form the Bochner identity gives
\[
 \|\nabla\omega\|_2^2+a_p\|\omega\|_2^2
 \le\int W_p|\omega|^2
 \le\|W_p\|_{n/2}\|\omega\|_{2n/(n-2)}^2.
\]
Apply (211.2) to $|\omega|$, using the weak Kato inequality
$|d|\omega||\le|\nabla\omega|$. It follows that
\[
 (1-b_p)\|\nabla\omega\|_2^2+
 (a_p-b_pV^{-2/n})\|\omega\|_2^2\le0.
\]
Both coefficients are strictly positive, proving the claim. The application
to the nonsmooth norm of a form follows by replacing it by
$(|\omega|^2+\delta^2)^{1/2}$ and passing to the limit; compactness of $M$
and smoothness of $\omega$ justify the passage. If these hypotheses hold in
every odd degree, the argument in part 1 proves positive Euler characteristic.
The quantities in (211.3) are compatible with metric rescaling: $a_p$ scales
as curvature, while the $L^{n/2}$ curvature norm is scale invariant.

\paragraph{Stress test.}
The matrix is \emph{not} a metric on a closed manifold with $\chi\le0$.
Realizing an algebraic tensor locally does not supply a global metric,
harmonic form, or topology. It only eliminates the pointwise Bochner
positivity lemma. Conversely, having some negative eigenvalues does not by
itself disprove the integral criterion; harmonic forms need not concentrate
on their negative eigenspaces.

Condition (211.3) is a genuine additional hypothesis. One may formally set
$W_p=(a_p-\lambda_{\min}(q_p))_+$, but positive sectional curvature has not
been shown here to make its norm small relative to the Sobolev constant.
Neither scaling nor compactness repairs this: scaling leaves the critical
norm unchanged, and compactness of the \emph{individual} manifold supplies
no uniform control over this class of metrics. At the equality threshold in
(211.3), the proof only gives nonnegativity and does not force vanishing.
In addition, vanishing of all odd cohomology could be unnecessarily strong;
a successful Euler-characteristic argument might instead control the
alternating sum without this vanishing. The present work does not establish
such a cancellation or an everywhere positive Euler density.

\paragraph{Outcome.}
\textbf{Outcome: Exploratory approach with unresolved gap.}

The attempt gives an explicit exact obstruction to a pointwise odd-form
Bochner proof, together with a fully proved integral criterion that would
replace that false premise. It recovers the elementary four-dimensional
case and pinpoints the first uncontrolled degree in dimension six. The
missing step is a global estimate of the negative curvature contribution,
relative to the Sobolev inequality, derived from sectional positivity alone,
or a different argument controlling Euler characteristic without odd-degree
vanishing. No counterexample to Hopf's conjecture, and no unrestricted proof,
is obtained.
% END RESEARCH ADDITION 211
\subsection{Sources}
\begin{itemize}\raggedright
\item[\textnormal{[S1]}]\hypertarget{source:211:1}{} Lee Kennard, Lawrence Mouille, Jan Nienhaus. On Hopf's conjecture and positive second intermediate Ricci curvature. 2025.
\url{https://arxiv.org/abs/2507.16936}.
{\small\textit{Catalogue formulation source.}}
% BEGIN RESEARCH REFERENCES 211
\item[\textnormal{[E1]}]\hypertarget{extra:211:1}{} Jan Nienhaus,
\emph{An improved four-periodicity theorem and a conjecture of Hopf with symmetry},
arXiv:2211.13151v1 (23 November 2022), abstract and main Hopf-with-$T^4$ theorem.
\url{https://arxiv.org/abs/2211.13151}.
\item[\textnormal{[E2]}]\hypertarget{extra:211:2}{} Uwe Semmelmann and Gregor Weingart,
\emph{The Weitzenb\"ock machine}, Compositio Mathematica 146 (2010), 507--540;
arXiv:math/0702031, Section 3, especially (3.9)--(3.10).
\url{https://arxiv.org/pdf/math/0702031}.
% END RESEARCH REFERENCES 211
\end{itemize}

\end{document}
